\section{Arguments and formalization}
\label{sec:levels-05}

\begin{definition}[Argument]
An \emph{argument} is an ordered pair \(\tuple{\Gamma,\varphi}\), where
\(\Gamma\) is a set of formulas called the premises and \(\varphi\) is a
formula called the conclusion.
\end{definition}

An argument \(\tuple{\Gamma,\varphi}\) is \emph{semantically valid} iff
\(\Gamma\entails\varphi\); for \(\Gamma=\varnothing\) this agrees with
validity of the formula \(\varphi\), and the phrase names no relation
beyond semantic consequence.  A proof of
the conclusion from the premises is a derivation witnessing
\(\Gamma\derives_{\mathcal C}\varphi\).  If every premise and the conclusion
are sentences, an invalid argument is refuted by a countermodel.  For formulas
with free variable occurrences, the refuting data are stated explicitly as a
structure and a variable assignment.

\begin{definition}[Countermodel]
Let every member of \(\Gamma\cup\set{\varphi}\) be a sentence.  A
\emph{countermodel} to \(\tuple{\Gamma,\varphi}\) is a model of \(\Gamma\)
that does not satisfy \(\varphi\).
The restriction to sentences preserves the fixed standing of a model as a
structure satisfying a specified set of sentences.
\end{definition}

\begin{method}[Formalizing an argument]
Perform the following steps in order.
\begin{enumerate}
  \item State the intended domain and the conditions under which an expression
        counts as a premise or conclusion.
  \item Declare every nonlogical symbol, including its arity and intended
        interpretation.
  \item Resolve pronouns, tense, modality, plurality, restricted quantification,
        and uniqueness before translating.
  \item Translate each premise separately.  Preserve quantifier scope with
        parentheses.
  \item Translate the conclusion separately.
  \item Test the translation against intended positive and negative cases.
  \item Only then determine derivability.  For a sentence argument, search for
        a countermodel when semantic consequence is in doubt; for formulas
        with free variable occurrences, search for an explicit structure and
        variable assignment that refute semantic consequence.
\end{enumerate}
\end{method}

\begin{example}[Quantifier order]
Over a domain of people, let \(A(x,y)\) mean that \(x\) admires \(y\).
The sentence
\[
  \forall x\,\exists y\,A(x,y)
\]
allows the admired person to depend on \(x\).  The sentence
\[
  \exists y\,\forall x\,A(x,y)
\]
requires one person admired by everyone.  Neither sentence may be substituted
for the other without changing the claim.
\end{example}

